\documentclass[11pt,a4paper]{article}

\usepackage{amsmath,amssymb,amsthm}
\usepackage{physics}
\usepackage{bm}
\usepackage{booktabs}
\usepackage{hyperref}
\usepackage{xcolor}
\usepackage{listings}
\usepackage[margin=2.5cm]{geometry}

\lstset{
  language=C++,
  basicstyle=\ttfamily\small,
  keywordstyle=\color{blue!70!black},
  commentstyle=\color{green!50!black},
  frame=single,
  breaklines=true,
}

\newcommand{\sgn}{\operatorname{sgn}}
\newcommand{\DT}{\Delta t}

\title{Stochastic Noise in the Generalised Phase Force:\\
       Mathematical Model and Simulation Implementation}
\author{CUFCM\_FIL}
\date{}

\begin{document}
\maketitle

% ============================================================
\section{Phase dynamics without noise}
% ============================================================

Each filament is characterised by a phase $\psi(t)\in[0,2\pi)$ that indexes
its position in the beat cycle.  In the dynamical (non-calibration) regime the
phase velocity $\dot\psi$ is not prescribed but is solved as part of the
force-balance problem.

\subsection{Generalised forces and the mobility system}

The simulation assembles a linear system coupling the filament degrees of
freedom (segment positions, phase, and---when enabled---beat-plane rotation
angle $\psi_2$) to the forces acting on them.  For the phase degree of freedom
the \emph{generalised phase force} is

\begin{equation}
  Q_\psi(\psi) = Q_{\text{ref}}(\psi)\,\frac{\omega_0}{2\pi},
  \label{eq:qphase_det}
\end{equation}

where $Q_{\text{ref}}(\psi)$ is a lookup table computed during the calibration
run (see Section~\ref{sec:calibration}) and $\omega_0$ is the filament's
natural angular frequency.  The factor $\omega_0/(2\pi)$ rescales the reference
force when a filament's natural frequency differs from the reference frequency.

The generalised force $Q_\psi$ enters the right-hand side of the mobility
linear system.  After solving, the phase velocity is obtained as

\begin{equation}
  \dot\psi = M_{\psi\psi}\,Q_\psi + (\text{hydrodynamic coupling terms}),
\end{equation}

where $M_{\psi\psi}$ is the $({\psi},{\psi})$ block of the mobility matrix.
The phase is then advanced by the forward-Euler step

\begin{equation}
  \psi(t+\DT) = \psi(t) + \dot\psi\,\DT.
  \label{eq:phase_det_update}
\end{equation}

% ============================================================
\section{Adding stochastic noise}
% ============================================================

Real cilia are not perfect oscillators: motor noise, thermal fluctuations, and
internal regulatory variability all perturb the beat rhythm.  We model this
with additive Gaussian white noise on the generalised phase force.

\subsection{Continuous stochastic equation}

The noisy generalised force is

\begin{equation}
  Q_\psi^{\text{stoch}}(t) = Q_{\text{ref}}(\psi)\,\frac{\omega_0}{2\pi}
    + \sigma\,\xi(t),
  \label{eq:qphase_stoch}
\end{equation}

where $\xi(t)$ is a unit Gaussian white-noise process,

\begin{equation}
  \langle\xi(t)\rangle = 0, \qquad
  \langle\xi(t)\,\xi(t')\rangle = \delta(t-t').
\end{equation}

The parameter $\sigma$ (runtime variable \texttt{FORCE\_NOISE\_MAG}) controls
the noise amplitude.  The resulting phase dynamics are governed by a
Langevin-type equation:

\begin{equation}
  \dot\psi = M_{\psi\psi}\!\left[Q_{\text{ref}}(\psi)\,\frac{\omega_0}{2\pi}
    + \sigma\,\xi(t)\right].
  \label{eq:langevin_phase}
\end{equation}

\subsection{Effective phase diffusion}

Separating the drift and noise contributions, the phase evolves as

\begin{equation}
  d\psi = \underbrace{M_{\psi\psi}\,Q_{\text{ref}}(\psi)\,
    \frac{\omega_0}{2\pi}\,dt}_{\text{deterministic drift}}
  + \underbrace{M_{\psi\psi}\,\sigma\,dW_t}_{\text{noise}},
\end{equation}

where $dW_t = \xi(t)\,dt$ is an increment of the Wiener process.  The
effective phase diffusion coefficient is therefore

\begin{equation}
  \boxed{D_{\psi} = \tfrac{1}{2}(M_{\psi\psi}\,\sigma)^2.}
  \label{eq:Deff}
\end{equation}

$D_\psi$ quantifies the rate of phase-variance growth for an isolated
filament:

\begin{equation}
  \operatorname{Var}[\psi(t) - \psi(0)] = 2\,D_\psi\,t.
  \label{eq:var_growth}
\end{equation}

\subsection{Euler--Maruyama discretisation}

Numerically the noise increment in force space is drawn at every timestep.
The discrete update is

\begin{equation}
  \psi_{n+1} = \psi_n + M_{\psi\psi}\!\left[Q_{\text{ref}}(\psi_n)\,
    \frac{\omega_0}{2\pi} + \sigma\,\xi_n\right]\DT,
  \qquad \xi_n \sim \mathcal{N}(0,1)\;\text{i.i.d.}
  \label{eq:EM_force}
\end{equation}

\textbf{Note on the $\sqrt{\DT}$ factor.}
Equation~\eqref{eq:EM_force} adds $\sigma\xi_n\DT$ to the force each step.
In phase space this corresponds to a phase kick of magnitude
$M_{\psi\psi}\sigma\DT$ per step.  Summing $T/\DT$ independent kicks gives

\begin{equation}
  \operatorname{Var}[\psi(T)] \approx
    (M_{\psi\psi}\sigma)^2 \frac{T}{\DT}\,(\DT)^2
  = (M_{\psi\psi}\sigma)^2\,\DT\,T.
\end{equation}

Comparing with equation~\eqref{eq:var_growth} yields $D_\psi = \tfrac{1}{2}
(M_{\psi\psi}\sigma)^2\DT$, which \emph{depends on the timestep}.  This
arises because the implementation applies the noise to the force (not directly
to the phase), so the correct Langevin scaling would require adding
$\sigma\sqrt{\DT}\,\xi_n$ instead of $\sigma\xi_n$.  The current code
(\texttt{mobility\_solver.cpp}, line~364) reads:

\begin{lstlisting}
q_phase += FORCE_NOISE_MAG*noise;   // noise ~ N(0,1), no sqrt(DT)
\end{lstlisting}

For a proper Langevin process, the line should be:

\begin{lstlisting}
q_phase += FORCE_NOISE_MAG*std::sqrt(DT)*noise;
\end{lstlisting}

With the corrected form $D_\psi = \tfrac{1}{2}(M_{\psi\psi}\sigma)^2$ becomes
truly timestep-independent.

% ============================================================
\section{Calibration and the reference force table}
\label{sec:calibration}
% ============================================================

The table $Q_{\text{ref}}(\psi)$ is produced by a separate calibration run
compiled with \texttt{WRITE\_GENERALISED\_FORCES=true}.  During calibration,
$\dot\psi = \omega_0$ is prescribed (not solved dynamically), so the
simulation simply measures the force required to maintain the reference beat
at each phase.  The calibration preset uses \texttt{force\_noise\_mag=0}
so that $Q_{\text{ref}}$ captures only the deterministic driving, keeping the
noise channel in the production run well-defined.

\begin{table}[h]
\centering
\begin{tabular}{lll}
\toprule
Flag & Calibration & Production \\
\midrule
\texttt{WRITE\_GENERALISED\_FORCES} & \texttt{true} & \texttt{false} \\
\texttt{DYNAMIC\_PHASE\_EVOLUTION}  & forced \texttt{false} & \texttt{true} \\
\texttt{force\_noise\_mag}          & $0$ & $\sigma$ (swept) \\
Active noise site & \texttt{filament.cpp:1392} (dead) & \texttt{mobility\_solver.cpp:364} \\
\bottomrule
\end{tabular}
\caption{Compile-time and runtime flag states for the two run modes.}
\end{table}

% ============================================================
\section{Effect of changing \texttt{FORCE\_NOISE\_MAG} ($\sigma$)}
% ============================================================

\subsection{Signal-to-noise ratio}

The deterministic driving force is calibrated to
$Q_{\text{ref}} \sim Q_0 \equiv 220$ (dimensionless).  The noise
perturbation per timestep is of order $\sigma$.  Define the signal-to-noise
ratio

\begin{equation}
  \mathrm{SNR} \equiv \frac{Q_0}{\sigma}.
\end{equation}

\begin{table}[h]
\centering
\begin{tabular}{ccc}
\toprule
Sim $i$ & $\sigma$ & $\mathrm{SNR}$ \\
\midrule
0 & $0.0$    & $\infty$ \\
1 & $0.22$   & $1000$ \\
2 & $0.46$   & $478$ \\
3 & $1.0$    & $220$ \\
4 & $2.15$   & $102$ \\
5 & $4.64$   & $47$ \\
6 & $10.0$   & $22$ \\
7 & $21.5$   & $10$ \\
8 & $46.4$   & $4.7$ \\
9 & $100.0$  & $2.2$ \\
\bottomrule
\end{tabular}
\caption{Log-uniform noise sweep in \texttt{pair\_noise\_sweep}.
$\mathrm{SNR} = 220/\sigma$; phase slips are expected when $\mathrm{SNR}
\lesssim 10$--$20$, i.e.\ sims 7--9.}
\end{table}

\subsection{Phase diffusion and slip between a pair}

For a pair of filaments (A and B) with equal natural frequencies
($\omega_A = \omega_B = \omega_0$, \texttt{pair\_dp=1}), the mean phase
difference

\begin{equation}
  \Delta\psi(t) \equiv \psi_A(t) - \psi_B(t)
\end{equation}

has zero deterministic drift.  Any long-term phase slip is purely
noise-driven.  Since the two filaments carry independent noise realisations,
the variance of $\Delta\psi$ grows as

\begin{equation}
  \operatorname{Var}[\Delta\psi(t)] = 2\,D_A\,t + 2\,D_B\,t = 4\,D_\psi\,t
  \quad (D_A = D_B = D_\psi),
\end{equation}

so the pair's relative phase diffusivity is $2D_\psi$.  A \emph{phase slip}
event occurs when $|\Delta\psi|$ exceeds $\pi$ (half a period) and the
filaments re-synchronise in the next cycle out of phase.

The slip rate measured in the simulation is

\begin{equation}
  \dot s \equiv \frac{\Delta\psi(t_\mathrm{end}) - \Delta\psi(t_\mathrm{start})}
                     {t_\mathrm{end} - t_\mathrm{start}}
  \quad [\text{cycles\,s}^{-1}],
\end{equation}

and scales as $\dot s \propto D_\psi \propto \sigma^2 M_{\psi\psi}^2$.

\subsection{Three dynamical regimes}

\paragraph{Weak noise ($\sigma \ll Q_0$, sims 0--4).}
$\Delta\psi$ performs a slow random walk around the synchronised state.
The pair remains approximately in phase over many beat cycles.  The slip rate
$\dot s \approx 0$.

\paragraph{Moderate noise ($\sigma \sim Q_0$, sims 5--7).}
Phase slips occur intermittently.  The pair alternates between epochs of
near-synchrony and epochs of relative sliding.  The slip rate grows rapidly
with $\sigma$.

\paragraph{Strong noise ($\sigma \gg Q_0$, sims 8--9).}
The noise overwhelms the deterministic driving.  The phase becomes
effectively diffusive ($\dot\psi \approx M_{\psi\psi}\sigma\xi$); the
filament beats erratically and the pair loses all phase coherence.

% ============================================================
\section{Summary of the parameter chain}
% ============================================================

\begin{equation*}
\underbrace{\sigma}_{\texttt{force\_noise\_mag}\text{ in }
  \texttt{globals.ini}}
\xrightarrow{\texttt{cilia\_sim\_main.cu:74}}
\underbrace{\texttt{FORCE\_NOISE\_MAG}}_{\text{runtime Real}}
\xrightarrow[\text{(production)}]{\texttt{mobility\_solver.cpp:364}}
Q_\psi^{\text{stoch}} = Q_{\text{ref}}\tfrac{\omega_0}{2\pi} + \sigma\xi_n
\xrightarrow{M_{\psi\psi}}
\dot\psi_n
\xrightarrow{\times\DT}
\psi_{n+1}
\end{equation*}

The single scalar $\sigma$ is thus the sole runtime control over the
stochastic intensity, requiring no recompilation.  Sweeping it
log-uniformly from $0$ to $100$ (with $Q_0 = 220$) spans the full range
from noise-free synchronisation to noise-dominated phase diffusion.

\end{document}
